\documentclass[11pt]{article}
\usepackage[margin=1in]{geometry}
\usepackage{amsmath,amssymb}
\usepackage[version=4]{mhchem}
\usepackage[T1]{fontenc}
\usepackage{lmodern}
\usepackage[authoryear,round]{natbib}
\usepackage{graphicx,booktabs,longtable,array,needspace}
\usepackage{url}
\usepackage[hidelinks]{hyperref}
\setlength{\emergencystretch}{2em}
\setlength{\tabcolsep}{4pt}
\setlength{\LTpre}{5pt}
\setlength{\LTpost}{6pt}
\newcolumntype{P}[1]{>{\raggedright\arraybackslash}p{#1}}
\newcommand{\paperfigure}[1]{\includegraphics[width=\linewidth,height=0.65\textheight,keepaspectratio]{figures/#1}}
\renewcommand{\thefigure}{S\arabic{figure}}
\renewcommand{\thetable}{S\arabic{table}}
\title{Supporting Information for\\Precession-phase dependence of abrupt warming occurrence in Greenland and speleothem records}
\author{}
\date{}
\begin{document}
\maketitle
\noindent This Supporting Information contains Texts S1--S7, Tables S1--S5, and Figures S1--S14. Machine-readable versions of the tables retain the precision, source identifiers, and diagnostics of the saved analyses. References to numbered main figures refer to the accompanying article.

\section*{Text S1. Event inventories, timing, and observation support}

The primary catalogue contains 34 Greenland Interstadial (GI) starts from the NGRIP event framework of \citet{rasmussen2014stratigraphic} and 21 transitions from two speleothem records: 16 from Melchsee--Frutt (MF; \citealp{fohlmeister2023role}) and five from Sofular Cave (\citealp{held2024dansgaard}). The underlying speleothem datasets are archived separately \citep{fohlmeister2023mfData,held2024sofularData}. NGRIP start ages are taken directly from the published framework; the source's original subevent labels and definition-error codes are retained in the processed table. Warming and cooling boundaries are jointly ordered for chronology sampling, but only warming starts enter the occurrence model. Table~\ref{tab:S1} lists the numerical inventory and source labels.

The speleothem inventory follows published event identities. Numerical times are assigned within local windows around literature anchors rather than by searching the entire record for new events. Each proxy is interpolated onto a 1-yr grid and Gaussian-smoothed with a nominal standard deviation of 50 yr. Interpolation and smoothing are performed independently within continuous segments separated by gaps exceeding 0.5 kyr. The event age is the position of the largest gradient in the expected transition direction within $\pm400$ yr of its anchor. Nine alternatives cross smoothing standard deviations of 35, 50, and 75 yr with half-windows of 300, 400, and 500 yr. These alternatives quantify sensitivity to the numerical timing rule; they are not nine independent event inventories or a calibrated probability distribution for proxy measurement error.

MF labels 6.1--6.16 retain their sequence numbering. Sofular labels 6.9--6.13 are renumbered as composite events 6.17--6.21. The visually assessed overlap (Figure~\ref{fig:S01}) places MF 6.16 near an unnumbered Sofular feature at approximately 175.2 kyr BP, whereas Sofular 6.9 is the next older published event. We therefore retain both MF 6.16 and Sofular 6.9 (Figure~\ref{fig:S02}). This comparison does not tune either chronology. The MF proxy gap extends from 174.9052 to 181.986 kyr BP; the selected MF sequence ends with event 6.16 at 174.752 kyr BP, and Sofular supplies the older continuation. Thus this 7.08-kyr gap in MF is not treated as a 7.08-kyr gap in the composite observation interval.

NGRIP observation support is 12--123 kyr BP. Its older endpoint follows the continuous, undisturbed record coverage described by \citet{ngrip2004high}; it is an analysis boundary rather than a detection limit stated by \citet{rasmussen2014stratigraphic}. The observed older tail contains no additional published GI start. Speleothem support is 132.5--204.5 kyr BP, informed by the coverage and published comparison in \citet{held2024dansgaard}. Because this interval extends beyond strict MIS~6, we call it a MIS~6-centered sequence; ``MIS6'' in file names is a project identifier. The oldest 1.5 kyr of each observation interval is reserved for event-history initialization, leaving response intervals of 12--121.5 and 132.5--203 kyr BP. Their combined response exposure is 180 kyr. The gap at 123--132.5 kyr BP contributes neither bins nor event history. The estimand is the warming-onset rate per unit of whole observation time, including observed event-free tails; it is not a rate conditional on an independently reconstructed cold-state duration. As in any analysis of a published event catalogue, interpretation assumes that the inventory is adequate over its declared support.

The Barker reconstruction is analyzed separately \citep{barker2011800}. We use its Supplementary Table~S3 SpeleoAge column and variable-threshold picks within 0--400 kyr, yielding 70 events. The 59 fixed-threshold picks are a subset of these events on the same source chronology and form a point-age definition sensitivity. Table~\ref{tab:S1} retains the source Excel row and fixed-threshold membership for every primary Barker event. The response interval is 0--398.5 kyr BP; 398.5--400 kyr supplies history only. All 70 variable-threshold and 59 fixed-threshold point ages fall within response support. Barker events are not pooled with NGRIP or the speleothem composite, and the two Barker definitions are not independent samples.

\section*{Text S2. Epochs and working chronology ensembles}

\subsection*{Age coordinates}

All analyses use kyr BP relative to 1950 CE. NGRIP ages in years before 2000 CE are converted as $t_{\rm BP}=t_{\rm b2k}/1000-0.05$; the same conversion applies to the absolute GICC05 time coordinate \citep{rasmussen2014stratigraphic,rasmussen2022gicc05,rasmussen2023ice}. MF archive units, Sofular source metadata, LR04 archive metadata, and the \ce{CO2} composite identify 1950-based ages, which require unit conversion only \citep{fohlmeister2023mfData,held2024sofularData,lisiecki2005pliocene,bereiter2015revision}. La2004's signed kyr coordinate is measured from J2000 and is negative in the past, giving $t_{\rm BP}=-t_{\rm source}-0.05$ \citep{laskar2004long}. The local solstice-insolation file has positive ages but lacks an explicit reference year; numerical reconstruction supports a J2000 coordinate, and 0.05 kyr is subtracted. Barker's exact SpeleoAge reference year remains unverified; its values are retained under the working BP1950 assumption. No epoch correction is applied to an error width or elapsed duration. Table~\ref{tab:S2} distinguishes verified epochs from inferences and assumptions.

\subsection*{NGRIP: correlated chronology and definition errors}

We generate 10,000 combined realizations of the 69 retained warming and cooling boundaries. Transition-definition codes are represented by independent zero-mean Gaussian errors, with working standard deviations of 20, 50, 200, 100, and 30 yr for codes a, b, c, d, and f, respectively; explicitly tabulated $\pm4$-yr values use a 4-yr standard deviation. These numerical representations of definition uncertainty are study assumptions attached to the published codes, not newly estimated sampling distributions.

For the layer-counted chronology, the annual GICC05 maximum counting error (MCE) curve determines a working standard deviation equal to MCE/2 \citep{svensson2008sixty,rasmussen2022gicc05}. The exact counted endpoint is 60.202 ka b2k (60.152 kyr BP), with MCE 2.611 kyr. Let $k_j$ denote the prepared age knots, spaced nominally 5 kyr apart with the counted endpoint explicitly retained, and let $V_j$ be the squared working chronology standard deviation at each knot. Starting with $V_{-1}=0$, draw
\begin{equation}
 e_j\sim N(0,V_j-V_{j-1}),\qquad \delta_j=\sum_{m\leq j}e_m.
\end{equation}
Thus the increment standard deviation is $\sqrt{V_j-V_{j-1}}$, not the difference of adjacent standard deviations. This construction yields cumulative, correlated errors while allowing local deformation. For an event between two knots, its chronology offset is a linear interpolation of the two offsets. Before conditioning, its variance at interpolation weight $w$ is $V_j+w^2(V_{j+1}-V_j)$, which need not equal the square of an interpolated envelope. The realized knot ages must increase strictly; the final event ages, after independent definition errors are added, must also retain the complete warming/cooling order. Whole proposals failing either condition are rejected without sorting events (Figure~\ref{fig:S03}).

Published age uncertainties for the flow-model extension remain unquantified \citep{corrick2020synchronous}. We adopt a nominal envelope equal to $\pm4.5\%$ of b2k age throughout the extension and interpret it approximately as $2\sigma$. \citet{moseley2020nalps19} extrapolated the counted age--uncertainty trend to approximately 4.5\% by 120 ka; that study does not prescribe a uniform percentage over the entire extension or the covariance model used here. Our envelope is therefore an explicit working scenario. The prepared grid preserves the counted endpoint and connects it continuously to the next modeled knot. The Gaussian envelope is not a hard bound, and the accepted quantiles are not calibrated chronology confidence intervals. The 2.5-, 5-, and 10-kyr alternatives in Text~S6 vary this working process. They do not reproduce a published age-model posterior or the distinct methodology of \citet{myrvollnilsen2022comprehensive}.

\subsection*{MF and Sofular: event timing and chronology}

For MF, the approximate 95\% age-model envelope supplies the chronology scale. Lower and upper bounds are ordered; if a nominal age lies outside the reported pair, the envelope is expanded to include it. This repair affects 188 source rows. Duplicate ages retain their outer envelope, and interpolation does not cross source gaps. At a detector-assigned event time $t$, the larger distance to the repaired bounds is divided by 1.95996 to obtain $\sigma(t)$. One standard-normal draw $z$ is shared by all MF events, giving $\delta(t)=z\sigma(t)$. This retains a simple record-wide dependence assumption; it is not a reconstruction of all age-model degrees of freedom.

Sofular uncertainties use the raw dated depths, corrected U--Th ages, reported $2\sigma$ errors, and published individual age--depth curves for So-4 and So-57 \citep{held2024dansgaard,held2024sofularData}. Write the reference curve as $g_0(d)$ and its age at dated depth $d_j$ as $m_j$. Independently draw $\epsilon_j\sim N(0,s_j^2)$, where $s_j$ is half the reported $2\sigma$ dating error, and set $u_j=m_j+\epsilon_j$. Between neighboring controls, define
\begin{equation}
 g_r(d)=(1-q)u_j+q u_{j+1},\qquad
 q=\frac{g_0(d)-m_j}{m_{j+1}-m_j}.
\end{equation}
This preserves the shape of the published reference growth curve between dated controls while perturbing its age coordinates. Nonmonotone control proposals are rejected. No interpolation crosses the So-4 hiatus at 781.5 mm or extrapolates beyond dated/proxy support. Ten So-4 controls in the relevant growth segment enter the main ensemble.

The point ages remain on the published stack. An event's detector age is projected onto an individual age--depth curve to locate the error field, and its realized stack age is $t+g_r(d_t)-g_0(d_t)$. This same-age projection is an assumed coordinate correspondence, not a verified physical event layer or reconstruction of stack-alignment transformations. Measured U--Th central ages and reference model ages need not coincide. The calculation transfers analytical dating errors onto a fixed reference curve; it does not rerun StalAge or iscam or reproduce their posteriors \citep{scholz2011stalage,fohlmeister2012statistical}.

Each realization chooses one of the nine timing settings independently for each source record, applies the record's chronology perturbation, and rejects any reversal in the complete composite event order. The main 10,000-member ensemble uses So-4 for all five Sofular events (Figure~\ref{fig:S04}). An alternative uses So-57 for source events 6.9--6.12 and So-4 for event 6.13, which lies outside the So-57 proxy support. Moving projection coordinates by $-1,-0.5,0,0.5,$ or 1 kyr evaluates dependence on the assumed transfer location; it changes error weights rather than point ages and is not a simulated physical lag. Proxy measurement noise, unknown between-date growth history, shared systematic dating errors, missing stack transformations, and inter-record climate-transmission lags are not propagated. Order conditioning can induce additional dependence, including at the MF--Sofular boundary.

\subsection*{Barker: control offsets on SpeleoAge}

Barker Supplementary Table~S1 provides 60 chronology controls with combined matching and absolute-dating errors \citep{barker2011800}. We use the printed combined values directly as uniform proposal half-widths $U_j$: $\delta_j\sim\mathrm{Uniform}(-U_j,U_j)$. The matching and dating components are not added again, and the combined values are not relabeled as Gaussian standard deviations. Proposals are independent before conditioning on strict control-age order. For an event of nominal SpeleoAge $t$ between controls $a_j,a_{j+1}$, its offset is
\begin{equation}
 \delta(t)=(1-w)\delta_j+w\delta_{j+1},\qquad
 w=\frac{t-a_j}{a_{j+1}-a_j}.
\end{equation}
All interpolation weights are calculated on SpeleoAge; no EDC3 coordinate is used in sampling. Before conditioning, the offset variance is $[(1-w)^2U_j^2+w^2U_{j+1}^2]/3$. The interpolated half-width is therefore an offset bound, not an interpolated standard deviation.

The long interval between controls at 264.24 and 317.70 kyr spans the approximately 265--315-kyr alignment gap and contains seven primary events. Their offsets use the same bracketing-control interpolation. Four events older than the last published control at 378.15 kyr are covered by an auxiliary control at 400 kyr. Its half-width is the larger of a linear extrapolation through the last four published uncertainty values and their maximum, yielding 1.30 kyr. That fixed amplitude is determined once; each proposal gives the auxiliary control its own independent uniform offset. The complete control map is rejected if its order reverses, and monotone interpolation then preserves event order automatically. Figure~\ref{fig:S05} shows the 10,000 accepted chronologies. This is a transparent working treatment of the gap and older tail, not a calibrated uncertainty reconstruction. It includes neither event-membership uncertainty nor additional picking error.

Source ensemble seeds are 20260907 (NGRIP), 20260904 (MF--Sofular), and 20260908 (Barker). Saved provenance records proposal counts, order rejections, input paths, and parameter settings. The original source files and nominal event ages are preserved.

\section*{Text S3. Conditional count model and null calibration}

Within each observation segment, let $y_i$ be the event count in a bin with center age $t_i$ and duration $\Delta t_i$. Nominal bin width is 0.2 kyr. Any shortened edge bin retains its actual duration in the offset, so $\mu_i=\Delta t_i\lambda_i$ and
\begin{equation}
 \ell=\sum_{i\in\mathcal R}\{y_i[\log\Delta t_i+\eta_i]-\Delta t_i\exp(\eta_i)-\log(y_i!)\},
 \qquad \eta_i=\log\lambda_i,
\end{equation}
where $\mathcal R$ denotes response bins. Conditional-count models can incorporate history while preserving a likelihood for sequential events \citep{davis2003observation,truccolo2005point}. Ages increase into the past; hence the preceding history of bin $i$ consists of strictly older bins. The covariate is
\begin{equation}
 H_i=\sum_{j:t_i<t_j\leq t_i+1.5\ {\rm kyr}}y_j,
\end{equation}
restricted to the same segment. Current-bin events do not enter their own history. This is a bin-center approximation, not exact continuous-event timing. The nominal 1.5-kyr window was motivated by the characteristic recurrence timescale of D--O variability and was not chosen by maximizing PI in the reported sensitivity grid.

LR04 and \ce{CO2} are interpolated at bin centers. Each scalar covariate is centered by its arithmetic mean over response bins and divided by its response-bin range; both source segments share the primary scaling. Primary reduced and full models are
\begin{align}
 \eta_i^{\rm red}&=\beta_0+\beta_H H_i+\beta_L L_i+\beta_C C_i+\beta_S S_i,\\
 \eta_i^{\rm full}&=\eta_i^{\rm red}+\beta_s\sin\phi_i+\beta_c\cos\phi_i.
\end{align}
Here $S_i=1$ for the speleothem segment and 0 for NGRIP. Barker uses the same model without $S_i$. The segment coefficient allows different baseline rates; it is not an age correction. No proxy-resolution covariate is included.

Phase is constructed from alternating local minima and maxima of the La2004 precession index \citep{laskar2004long}. A minimum has phase 0, a maximum phase $\pi$, and unwrapped phase increases linearly between successive extrema toward older ages. Response-bin phase is evaluated at bin centers; descriptive event phases are evaluated at event ages. Define
\begin{equation}
 A=(\beta_s^2+\beta_c^2)^{1/2},\qquad
 \phi_{\max}=\operatorname{atan2}(\beta_s,\beta_c)\pmod{2\pi},\qquad
 R_{\max/\min}=\exp(2A).
\end{equation}
The plotted phase multiplier is $\exp(\beta_s\sin\phi+\beta_c\cos\phi)$ at fixed background and history, rather than the total event rate. If both coefficients are zero, phase is unidentified. Rayleigh statistics provide an unconditional description of event-phase concentration \citep{mardia2000directional}; they test a different hypothesis from the conditional phase addition.

PI is $(\ell_{\rm full}-\ell_{\rm red})/(N\ln2)$, where $N$ is the response-event count, and $LR=2(\ell_{\rm full}-\ell_{\rm red})$. PI quantifies improvement in fitted likelihood per event; no held-out forecast skill is established. Nominal phase-test p values use a $\chi^2_2$ reference. AICc is calculated as $-2\ell+2k+2k(k+1)/(n-k-1)$ with $n$ equal to the number of response bins, consistent with the implemented count likelihood. It is used as a supplementary model comparison on identical support. We do not treat the response-bin count as the number of independent climate events.

Maximum-likelihood fitting uses L-BFGS-B with an analytic gradient, a maximum of 4,000 iterations, and objective tolerance $10^{-11}$ in SciPy \citep{virtanen2020scipy}. The intercept has numerical bounds $[-20,5]$ and other coefficients $[-20,20]$; the linear predictor has safety clipping limits $[-50,20]$. Saved diagnostics check finite fits, optimizer convergence, likelihood nesting, and whether clipping occurs. All reported main, valid-age, and bootstrap fits pass these checks. Alternative-specification outputs also record design rank, condition numbers, and coefficient-bound diagnostics. These are numerical safeguards, not evidence that the statistical model is correctly specified.

For each primary point-age fit, 9,999 catalogues are generated under its fitted reduced model \citep{davison1997bootstrap}. Simulation proceeds from the oldest to the youngest bin, including the history buffer; histories are updated from simulated counts and reset at segment boundaries. History before the oldest observation edge is initialized to zero. Both reduced and full models are refitted to each simulated catalogue, with all nuisance coefficients reestimated. Event totals are random rather than fixed at their observed counts. The empirical p value is
\begin{equation}
 p_{\rm boot}=\frac{1+\sum_{b=1}^{B}\mathbf1(LR_b^*\geq LR_{\rm obs})}{B+1}.
\end{equation}
The primary catalogue has 21 exceedances and $p_{\rm boot}=0.0022$ (Figure~\ref{fig:S08}); Barker has 137 and $p_{\rm boot}=0.0138$ (Figure~\ref{fig:S09}). No replicate requires rejection or retry. Seeds are 20260905 and 20260909, respectively. Exact binomial intervals in Table~\ref{tab:S3} describe finite-simulation precision of the unadjusted exceedance probability. This calibration uses fixed forcing grids and point-age fits; it is separate from age MC and is not assigned to alternative null hypotheses.

\section*{Text S4. Chronology sensitivity of the conditional results}

Independent random permutations pair the 10,000 NGRIP and 10,000 speleothem realizations without replacement (seed 20260906), using each source row exactly once. Every realization retains event identities, recomputes bin counts and histories, and refits both models. Forcing grids, climate scaling, observation support, and the 180-kyr response exposure remain fixed. Source events outside observation support cause a realization to be flagged, rather than clipped, resampled, or assigned a zero effect. Events remaining within observation support but entering a history-only buffer contribute history while leaving the response count.

Of 10,000 paired primary realizations, 18 are outside observation support and 9,982 have valid fits. Among valid fits, 66 contain 54 response events and 9,916 contain 55. Median PI is 0.16215 bits/event, with 2.5th--97.5th percentiles of 0.08616--0.21746; preferred phase is $330.04^\circ$ ($315.89$--$343.66^\circ$), and the phase rate ratio is 4.53 (2.92--5.96). Nominal p is below 0.05 in 9,851/9,982 valid fits (98.69\%). Across all source draws, at least 98.51\% are known to satisfy that threshold; no result is inferred for the 18 unsupported draws. Full-minus-reduced AICc is negative in 9,971 valid fits (Figure~\ref{fig:S06}; Table~\ref{tab:S3}).

All 10,000 Barker chronologies are supported and retain 70 response events. Median PI is 0.07763 bits/event (0.05350--0.10610), preferred phase $340.46^\circ$ ($331.27$--$351.67^\circ$), and rate ratio 2.66 (2.23--3.18). Nominal p is below 0.05 in 8,966 draws (89.66\%; Figure~\ref{fig:S07}). Each curve band is computed pointwise over phase. These quantiles and threshold fractions summarize sensitivity to specified chronology perturbations; they are neither process-sampling confidence intervals nor empirical null p values.

\section*{Text S5. Effect precision under full-model simulations}

Three analyses separate effect precision from the reduced-model test of association. A uses the 9,982 valid chronology refits described above. B simulates 5,000 catalogues from the point-age fitted full model. C samples 200 distinct valid chronology fits uniformly and generates 50 full-model catalogues from each, yielding 10,000 refits. Of the selected C generators, 199 contain 55 response events and one contains 54. Full-model simulation uses the same older-to-younger dynamic histories, zero initial history at each segment boundary, actual bin durations, and fixed external forcing grids as the null simulation. Each simulated catalogue is refitted for all coefficients, and counts may vary. No additional age perturbation is applied to simulated events in C: chronology uncertainty enters through the fitted generator. Seed 20260908 determines generator selection and simulations. All 15,000 refits pass numerical checks without retries.

For B, write $\widehat\theta=(\widehat\beta_s,\widehat\beta_c)'$ and let $C$ be the covariance of full-model bootstrap coefficient estimates $\widehat\theta_b^*$. The 95th percentile of
\begin{equation}
 Q_b=(\widehat\theta_b^*-\widehat\theta)'C^{-1}(\widehat\theta_b^*-\widehat\theta)
\end{equation}
is $q_{0.95}=6.27709$. An approximate joint coefficient confidence region is
\begin{equation}
 \{\theta:(\widehat\theta-\theta)'C^{-1}(\widehat\theta-\theta)\leq q_{0.95}\}.
\end{equation}
Calibration is referenced to the generator rather than the bootstrap mean, retaining finite-sample bias in $Q_b$. Projecting this region onto phase, rate ratio, and phase-response curves yields marginal parameter ranges and a simultaneous curve envelope. This is a plug-in bootstrap approximation under the fitted event process, not an exact finite-sample coverage guarantee. The current region excludes the zero-effect coefficient pair; a region containing zero would not identify a preferred phase.

The B phase projection is $286.53$--$374.08^\circ$, and the rate-ratio projection is 1.56--15.58. C reports central 95\% raw-refit working ranges: $295.92$--$365.45^\circ$ and 2.10--14.19, with medians $329.22^\circ$ and 5.09. Angles are unwrapped around the point estimate, $330.00^\circ$; an upper endpoint above $360^\circ$ continues into the next cycle. C retains estimator bias and is neither a calibrated combined confidence interval nor a posterior. B and C intervals need not nest because their constructions differ. For comparison using the same raw-refit quantile definition, the B phase range is $301.17$--$358.66^\circ$, narrower than C. Main Figure~3 and Table~\ref{tab:S3} distinguish these definitions. No corresponding full-model effect-precision study was performed for Barker, and its age range must not be substituted for sampling uncertainty.

\section*{Text S6. Design, background shape, memory, and event definition}

A complete 30-cell design experiment crosses history windows of 1, 1.5, 2, 3, and 5 kyr; bin widths of 0.1, 0.2, and 0.5 kyr; and origins shifted by zero or half a bin. All designs use the common response core 12--118 kyr BP for NGRIP and 132.5--199.5 kyr BP for speleothems, totaling 173 kyr and retaining the same 55 events. PI ranges from 0.13009 to 0.18908 bits/event, nominal p from 0.000740 to 0.007016, and phase from $323.82$ to $331.39^\circ$; full-minus-reduced AICc is negative in every design. At the primary bin/history settings, the common core gives PI 0.17249 versus 0.17933 on the 180-kyr main support (Figure~\ref{fig:S10}). Table~\ref{tab:S4} reports all cells, avoiding selection of a favorable setting.

Background-shape alternatives add LR04 squared, \ce{CO2} squared, both squares, or both squares plus their product, always retaining first-order terms. Transformations use the original scaled predictors, and each background specification appears in both reduced and full models; the nested addition remains the two phase terms. The five models share 180-kyr support and 55 events. PI remains 0.17051--0.17982 bits/event, phase $322.25$--$330.00^\circ$, and rate ratio 4.70--5.01. Adding LR04 squared lowers full-model AICc by 3.94 relative to the linear-background full model and shifts the peak toward $322.25^\circ$, demonstrating sensitivity of peak location to climate-response shape while retaining a phase likelihood gain.

Memory alternatives replace the prior-window count by elapsed time $\tau_i$ since the nearest strictly older occupied bin, or by $\log(1+\tau_i/1\ {\rm kyr})$, each with one coefficient. They retain the bin-center convention and update history only in younger bins. Prior waiting time is unknown before each segment's oldest observed event. We condition on its occupied bin and score only younger complete bins: GI-25 at 115.320 kyr BP lies in 115.2--115.4 kyr BP, and composite event 6.21 at 194.238 kyr BP lies in 194.1--194.3 kyr BP. Those events initialize later histories but do not enter the response likelihood. The resulting matched support is 12--115.2 and 132.5--194.1 kyr BP, totaling 164.8 kyr, 824 bins, and 53 events. The original count model is also refitted on this support. Its PI is 0.19612 bits/event, compared with 0.18832 for elapsed time and 0.19195 for log elapsed time (Figure~\ref{fig:S11}). This comparison changes the memory representation without confounding it with a different support; it is still a binned conditional-count model rather than an exact renewal-process likelihood.

A segment-pooling diagnostic adds segment-by-sine and segment-by-cosine terms to the common-phase full model. The added LR is 0.4276, nominal $p=0.8075$, and AICc increases by 3.65. Fitted phase maxima are $331.20^\circ$ for NGRIP and $327.43^\circ$ for speleothems, with rate ratios 4.07 and 6.95. The available data do not require separate phase coefficients; the nonsignificant interaction does not establish identical responses.

Barker's 59 fixed-threshold events yield PI 0.10423 bits/event, nominal $p=0.01409$, preferred phase $350.33^\circ$, and rate ratio 3.10 on the same support as its 70-event primary definition (Figure~\ref{fig:S12}). Counts and histories are recomputed for each definition. Bootstrap calibration, age MC, orbital alternatives, and LR04 interactions apply only to the variable-threshold definition; no corresponding fixed-threshold uncertainty is inferred.

For chronology structure, the complete So-57 overlap alternative gives median primary PI 0.16149 bits/event and phase $330.13^\circ$, with 98.53\% of 9,982 valid fits having nominal $p<0.05$. A separate 2,000-draw screening for each NGRIP knot spacing of 2.5, 5, and 10 kyr gives valid-draw fractions of 98.29\%, 98.40\%, and 99.00\%, respectively. The same paired speleothem ages are held fixed across knot spacings, but changing the NGRIP innovation dimension changes its error-process basis and between-knot marginal scales. Projection-position diagnostics for Sofular are also retained in Table~\ref{tab:S4}. These alternatives assess working assumptions; none supplies a calibrated chronology covariance. All specification tests use nominal likelihood-ratio references and do not inherit the main bootstrap p value.

\section*{Text S7. Orbital alternatives and LR04 modulation}

We compare eight models per catalogue: background only (BG), BG plus phase (Pre), and BG plus each of eccentricity, obliquity, or 65$^\circ$N summer-solstice daily mean top-of-atmosphere insolation, with and without Pre. BG includes the original climate and history terms and, for the primary catalogue, its segment term. Eccentricity is dimensionless, obliquity is converted from radians to degrees, and insolation is in W m$^{-2}$ \citep{laskar2004long}. Scalar orbital predictors receive the same response-bin mean/range scaling as the background covariates. The local insolation field uses solar longitude $90^\circ$ and solar constant 1,365 W m$^{-2}$. It is sampled natively every 1 kyr; linear interpolation onto model bins does not create new 0.2-kyr orbital information. Its epoch audit uses agreement with the corresponding La2004 numerical reconstruction (Table~\ref{tab:S2}).

There are ten nested comparisons: the original phase reference, then, for each scalar orbital variable, its addition to BG, its addition after Pre, and addition of Pre after that variable. A scalar adds one coefficient; phase adds two. Nine additional tests form a separate Holm family within each catalogue \citep{holm1979simple}; the original phase reference is excluded from that family. Table~\ref{tab:S5} gives all nominal and Holm-adjusted p values, AIC/AICc differences, and chronology ranges. Phase still adds PI after eccentricity or obliquity; after insolation its gain is 0.07870 bits/event in the primary catalogue and 0.03830 in Barker. Shared likelihood information is expected because these orbital descriptions covary; these comparisons do not isolate a unique causal driver or exclude indirect effects absorbed by LR04 and \ce{CO2}. Eccentricity enters here as an additive predictor, not an eccentricity-by-phase amplitude interaction.

For age sensitivity, the orbital and LR04-interaction analyses reuse the same 500 selected source realizations within each catalogue (selection seed 20260909). One selected primary realization is unsupported, leaving 499 valid comparisons; all 500 Barker draws are valid. All models within a realization share events, bins, forcing values, and exposure. The positive PI of a nested addition is not by itself a significance criterion, and these age ranges do not replace a sampling-based null calibration.

The LR04-modulation experiment adds $L_i\sin\phi_i$ and $L_i\cos\phi_i$ to the existing full model, retaining all main effects. At an LR04 value $L$, phase coefficients become $\beta_s+\gamma_s L$ and $\beta_c+\gamma_c L$. Displayed curves evaluate this one continuous fitted model at the exposure-weighted 25th, 50th, and 75th percentiles of LR04; they are not independently fitted climate groups. The primary interaction adds PI 0.02585 bits/event, LR 1.9709, nominal $p=0.3733$, and $\Delta\mathrm{AICc}=+2.11$ (Figure~\ref{fig:S13}). Barker adds 0.05613 bits/event, LR 5.4464, nominal $p=0.06566$, and $\Delta\mathrm{AICc}=-1.42$ (Figure~\ref{fig:S14}). Neither point-age nominal test reaches 0.05. The curve separation and chronology-PI distributions therefore do not establish LR04 modulation, although they describe an exploratory hypothesis for better-powered data.

Finally, chronology perturbations do not remove structural dependence among age models and orbital information. LR04 is orbitally tuned; the modeled Greenland extension inherits the ss09sea framework, whose development used a marine-isotope-stage age constraint derived from SPECMAP; and Barker's reconstruction is matched to speleothem ages \citep{lisiecki2005pliocene,wolff2010millennial,southon2004radiocarbon,barker2011800}. The Greenland chronology statement does not imply that individual events were tuned to precession. These dependencies limit the independence of the evidence even when perturbed event orders are preserved. The analyses quantify conditional associations under explicit chronology scenarios; they do not demonstrate orbital causation independently of chronology construction.

\clearpage
\setcounter{table}{0}\refstepcounter{table}\label{tab:S1}
\section*{Table S1. Event inventory and exposure}
% Generated by paper_tables.py; edit the exporter, not this file.
\Needspace{8\baselineskip}
\subsection*{(a) Nominal observation and response support (kyr BP)}
{\small
\begin{longtable}{llllll}
\toprule
Segment & Obs. start & Obs. end & Resp. start & Resp. end & Exposure \\
\midrule
\endfirsthead
\toprule
Segment & Obs. start & Obs. end & Resp. start & Resp. end & Exposure \\
\midrule
\endhead
\midrule
\multicolumn{6}{r}{Continued on next page} \\
\midrule
\endfoot
\bottomrule
\endlastfoot
NGRIP & 12.0 & 123.0 & 12.0 & 121.5 & 109.5 \\
MIS6 & 132.5 & 204.5 & 132.5 & 203.0 & 70.5 \\
Barker & 0.0 & 400.0 & 0.0 & 398.5 & 398.5 \\
\end{longtable}
}
Exposure is in kyr; ages increase into the past. The combined primary exposure is 180 kyr.
\Needspace{8\baselineskip}
\subsection*{(b) Primary published-event inventory}
{\small
\begin{longtable}{lllr}
\toprule
Event & Record & Source label & Age (kyr BP) \\
\midrule
\endfirsthead
\toprule
Event & Record & Source label & Age (kyr BP) \\
\midrule
\endhead
\midrule
\multicolumn{4}{r}{Continued on next page} \\
\midrule
\endfoot
\bottomrule
\endlastfoot
GI-1 & NGRIP & GI-1e & 14.642 \\
GI-2.1 & NGRIP & GI-2.1 & 22.970 \\
GI-2.2 & NGRIP & GI-2.2 & 23.290 \\
GI-3 & NGRIP & GI-3 & 27.730 \\
GI-4 & NGRIP & GI-4 & 28.850 \\
GI-5.1 & NGRIP & GI-5.1 & 30.790 \\
GI-5.2 & NGRIP & GI-5.2 & 32.450 \\
GI-6 & NGRIP & GI-6 & 33.690 \\
GI-7 & NGRIP & GI-7c & 35.430 \\
GI-8 & NGRIP & GI-8c & 38.170 \\
GI-9 & NGRIP & GI-9 & 40.110 \\
GI-10 & NGRIP & GI-10 & 41.410 \\
GI-11 & NGRIP & GI-11 & 43.290 \\
GI-12 & NGRIP & GI-12c & 46.810 \\
GI-13 & NGRIP & GI-13c & 49.230 \\
GI-14 & NGRIP & GI-14e & 54.170 \\
GI-15.1 & NGRIP & GI-15.1 & 54.950 \\
GI-15.2 & NGRIP & GI-15.2 & 55.750 \\
GI-16.1 & NGRIP & GI-16.1c & 57.990 \\
GI-16.2 & NGRIP & GI-16.2 & 58.230 \\
GI-17.1 & NGRIP & GI-17.1c & 59.030 \\
GI-17.2 & NGRIP & GI-17.2 & 59.390 \\
GI-18 & NGRIP & GI-18 & 64.050 \\
GI-19.1 & NGRIP & GI-19.1 & 69.570 \\
GI-19.2 & NGRIP & GI-19.2 & 72.290 \\
GI-20 & NGRIP & GI-20c & 76.390 \\
GI-21.1 & NGRIP & GI-21.1e & 84.710 \\
GI-21.2 & NGRIP & GI-21.2 & 85.010 \\
GI-22 & NGRIP & GI-22g & 89.990 \\
GI-23.1 & NGRIP & GI-23.1 & 103.990 \\
GI-23.2 & NGRIP & GI-23.2 & 104.470 \\
GI-24.1 & NGRIP & GI-24.1c & 106.700 \\
GI-24.2 & NGRIP & GI-24.2 & 108.230 \\
GI-25 & NGRIP & GI-25c & 115.320 \\
MIS 6.1 & MF & MF 6.1 & 134.976 \\
MIS 6.2 & MF & MF 6.2 & 138.393 \\
MIS 6.3 & MF & MF 6.3 & 146.251 \\
MIS 6.4 & MF & MF 6.4 & 146.847 \\
MIS 6.5 & MF & MF 6.5 & 148.439 \\
MIS 6.6 & MF & MF 6.6 & 149.458 \\
MIS 6.7 & MF & MF 6.7 & 150.462 \\
MIS 6.8 & MF & MF 6.8 & 151.115 \\
MIS 6.9 & MF & MF 6.9 & 159.522 \\
MIS 6.10 & MF & MF 6.10 & 160.879 \\
MIS 6.11 & MF & MF 6.11 & 165.006 \\
MIS 6.12 & MF & MF 6.12 & 168.566 \\
MIS 6.13 & MF & MF 6.13 & 169.675 \\
MIS 6.14 & MF & MF 6.14 & 171.506 \\
MIS 6.15 & MF & MF 6.15 & 172.364 \\
MIS 6.16 & MF & MF 6.16 & 174.752 \\
MIS 6.17 & Sofular & D-O 6.9 & 176.779 \\
MIS 6.18 & Sofular & D-O 6.10 & 179.831 \\
MIS 6.19 & Sofular & D-O 6.11 & 180.309 \\
MIS 6.20 & Sofular & D-O 6.12 & 190.972 \\
MIS 6.21 & Sofular & D-O 6.13 & 194.238 \\
\end{longtable}
}
NGRIP ages are published GI starts; MF and Sofular ages are maximum-gradient assignments near published labels (Text S1).
\Needspace{8\baselineskip}
\subsection*{(c) Barker variable-threshold inventory and fixed-threshold membership}
{\small
\begin{longtable}{lrl}
\toprule
Source identifier & Age (kyr BP) & Fixed threshold \\
\midrule
\endfirsthead
\toprule
Source identifier & Age (kyr BP) & Fixed threshold \\
\midrule
\endhead
\midrule
\multicolumn{3}{r}{Continued on next page} \\
\midrule
\endfoot
\bottomrule
\endlastfoot
Barker S3 row 10 & 11.402 & Yes \\
Barker S3 row 11 & 14.573 & Yes \\
Barker S3 row 12 & 27.979 & No \\
Barker S3 row 13 & 29.831 & Yes \\
Barker S3 row 14 & 32.920 & Yes \\
Barker S3 row 15 & 33.893 & Yes \\
Barker S3 row 16 & 35.398 & Yes \\
Barker S3 row 17 & 37.459 & Yes \\
Barker S3 row 18 & 38.566 & Yes \\
Barker S3 row 19 & 41.700 & Yes \\
Barker S3 row 20 & 43.694 & Yes \\
Barker S3 row 21 & 47.593 & Yes \\
Barker S3 row 22 & 53.867 & Yes \\
Barker S3 row 23 & 57.659 & Yes \\
Barker S3 row 24 & 59.302 & Yes \\
Barker S3 row 25 & 64.654 & Yes \\
Barker S3 row 26 & 71.615 & Yes \\
Barker S3 row 27 & 75.816 & Yes \\
Barker S3 row 28 & 83.583 & No \\
Barker S3 row 29 & 90.666 & Yes \\
Barker S3 row 30 & 103.355 & No \\
Barker S3 row 31 & 109.778 & Yes \\
Barker S3 row 32 & 125.328 & Yes \\
Barker S3 row 33 & 129.075 & Yes \\
Barker S3 row 34 & 147.366 & Yes \\
Barker S3 row 35 & 149.083 & No \\
Barker S3 row 36 & 150.972 & Yes \\
Barker S3 row 37 & 160.694 & Yes \\
Barker S3 row 38 & 164.034 & Yes \\
Barker S3 row 39 & 169.484 & Yes \\
Barker S3 row 40 & 173.183 & Yes \\
Barker S3 row 41 & 176.295 & Yes \\
Barker S3 row 42 & 177.959 & Yes \\
Barker S3 row 43 & 178.650 & Yes \\
Barker S3 row 44 & 189.979 & Yes \\
Barker S3 row 45 & 192.679 & Yes \\
Barker S3 row 46 & 199.077 & Yes \\
Barker S3 row 47 & 216.848 & Yes \\
Barker S3 row 48 & 226.714 & Yes \\
Barker S3 row 49 & 229.582 & Yes \\
Barker S3 row 50 & 232.027 & No \\
Barker S3 row 51 & 242.452 & Yes \\
Barker S3 row 52 & 243.886 & Yes \\
Barker S3 row 53 & 248.509 & No \\
Barker S3 row 54 & 251.876 & No \\
Barker S3 row 55 & 257.645 & No \\
Barker S3 row 56 & 259.518 & Yes \\
Barker S3 row 57 & 261.985 & Yes \\
Barker S3 row 58 & 264.305 & Yes \\
Barker S3 row 59 & 276.909 & Yes \\
Barker S3 row 60 & 278.806 & Yes \\
Barker S3 row 61 & 281.477 & Yes \\
Barker S3 row 62 & 291.430 & Yes \\
Barker S3 row 63 & 297.528 & Yes \\
Barker S3 row 64 & 301.831 & Yes \\
Barker S3 row 65 & 306.889 & Yes \\
Barker S3 row 66 & 317.780 & Yes \\
Barker S3 row 67 & 327.901 & Yes \\
Barker S3 row 68 & 334.861 & Yes \\
Barker S3 row 69 & 335.210 & Yes \\
Barker S3 row 70 & 335.744 & No \\
Barker S3 row 71 & 336.253 & Yes \\
Barker S3 row 72 & 349.383 & No \\
Barker S3 row 73 & 351.123 & Yes \\
Barker S3 row 74 & 366.805 & Yes \\
Barker S3 row 75 & 372.804 & Yes \\
Barker S3 row 76 & 378.191 & Yes \\
Barker S3 row 77 & 387.378 & Yes \\
Barker S3 row 78 & 392.246 & Yes \\
Barker S3 row 79 & 396.464 & No \\
\end{longtable}
}

\clearpage
\setcounter{table}{1}\refstepcounter{table}\label{tab:S2}
\section*{Table S2. Age epochs and uncertainty assumptions}
% Generated by paper_tables.py; edit the exporter, not this file.
\Needspace{8\baselineskip}
\subsection*{(a) Epoch harmonization}
{\small
\begin{longtable}{P{3.0cm}P{2.0cm}P{3.0cm}P{6.0cm}}
\toprule
Source & Epoch & Conversion & Evidence \\
\midrule
\endfirsthead
\toprule
Source & Epoch & Conversion & Evidence \\
\midrule
\endhead
\midrule
\multicolumn{4}{r}{Continued on next page} \\
\midrule
\endfoot
\bottomrule
\endlastfoot
NGRIP starts and GICC05 & 2000 CE & $t_{\rm b2k}({\rm yr})/1000-0.05$ & Explicit source definition \\
MF stack & 1950 CE & kyr ages unchanged & Archive calendar-BP unit \\
Sofular stack / So-4 & 1950 CE & kyr ages unchanged & Explicit article / raw headers \\
So-57 U-Th controls & 1950 CE & yr ages / 1000 & Explicit raw header \\
LR04 & 1950 CE & kyr ages unchanged & Archive and values checked \\
CO$_2$ composite & 1950 CE & yr ages / 1000 & Archive and values checked \\
La2004 precession, eccentricity, obliquity & J2000 & -signed kyr - 0.05 & Official epoch and values checked \\
65$^\circ$N solstice insolation & J2000 inferred & positive kyr - 0.05 & Numerical reconstruction; metadata not explicit \\
Barker SpeleoAge & 1950 CE assumed & kyr ages unchanged & Exact source epoch unverified \\
\end{longtable}
}
\Needspace{8\baselineskip}
\subsection*{(b) Working uncertainty assumptions}
{\small
\begin{longtable}{P{2.1cm}P{3.5cm}P{4.5cm}P{3.9cm}}
\toprule
Source & Error scale & Dependence & Boundary rule \\
\midrule
\endfirsthead
\toprule
Source & Error scale & Dependence & Boundary rule \\
\midrule
\endhead
\midrule
\multicolumn{4}{r}{Continued on next page} \\
\midrule
\endfoot
\bottomrule
\endlastfoot
NGRIP definition & Published location codes; working Gaussian SDs & Independent event errors: a 20, b 50, c 200, d 100, f 30 yr; explicit 4 yr retained & Reject complete crossed warming/cooling sequence \\
NGRIP counted & MCE / 2 as working SD & Cumulative independent variance increments at age knots; interpolate offsets & Reject nonmonotone knot map \\
NGRIP extension & 4.5\% of b2k age as approximate 2 SD envelope & Same cumulative process; preserve counted endpoint & Envelope is not a hard bound or published posterior \\
MF chronology & Outer approximate 95\% envelope / 1.95996 & One Gaussian factor shared by all MF events & Repair envelope to include point age; no interpolation across source gap \\
Sofular chronology & Reported U-Th 2 SD errors / 2 & Independent dated-depth errors transferred through reference growth curve & Monotone controls; no hiatus crossing or extrapolation \\
Speleothem timing & Nine equally weighted parameter settings & One smoothing/window choice per source record & Fixed event identity and final order \\
Barker chronology & Published combined errors as uniform half-widths & Independent control proposals; linear offset interpolation on SpeleoAge & Reject crossed controls; gap interpolation and auxiliary endpoint \\
\end{longtable}
}
Error widths are durations and receive unit conversion, never an epoch shift. Full source evidence and file paths are supplied in the accompanying CSVs.

\clearpage
\setcounter{table}{2}\refstepcounter{table}\label{tab:S3}
\section*{Table S3. Main estimates and uncertainty definitions}
% Generated by paper_tables.py; edit the exporter, not this file.
\Needspace{8\baselineskip}
\subsection*{(a) Point-age estimates}
{\small
\begin{longtable}{lrr}
\toprule
Quantity & Primary & Barker \\
\midrule
\endfirsthead
\toprule
Quantity & Primary & Barker \\
\midrule
\endhead
\midrule
\multicolumn{3}{r}{Continued on next page} \\
\midrule
\endfoot
\bottomrule
\endlastfoot
Response events & 55 & 70 \\
Exposure (kyr) & 180 & 398.5 \\
PI (bits/event) & 0.17933 & 0.08936 \\
Likelihood ratio & 13.673 & 8.6716 \\
Nominal p & 0.0010737 & 0.013091 \\
Preferred phase (degrees) & 330 & 336.91 \\
Maximum/minimum rate ratio & 4.938 & 2.8848 \\
Full-minus-reduced AICc & -9.6148 & -4.6494 \\
Descriptive Rayleigh p & 0.057959 & 0.079489 \\
\end{longtable}
}
\Needspace{8\baselineskip}
\subsection*{(b) Reduced-model null calibration}
{\small
\begin{longtable}{lrrlll}
\toprule
Catalogue & Simulations & Exceedances & Bootstrap p & 95\% lower & 95\% upper \\
\midrule
\endfirsthead
\toprule
Catalogue & Simulations & Exceedances & Bootstrap p & 95\% lower & 95\% upper \\
\midrule
\endhead
\midrule
\multicolumn{6}{r}{Continued on next page} \\
\midrule
\endfoot
\bottomrule
\endlastfoot
Primary & 9999 & 21 & 0.00220 & 0.00130 & 0.00321 \\
Barker & 9999 & 137 & 0.01380 & 0.01152 & 0.01618 \\
\end{longtable}
}
The interval is an exact binomial interval for the unadjusted null exceedance probability; the reported p uses the plus-one correction.
\Needspace{8\baselineskip}
\subsection*{(c) Chronology sensitivity}
{\small
\begin{longtable}{llrl}
\toprule
Catalogue & Quantity & Median [2.5\%, 97.5\%] & Valid/total \\
\midrule
\endfirsthead
\toprule
Catalogue & Quantity & Median [2.5\%, 97.5\%] & Valid/total \\
\midrule
\endhead
\midrule
\multicolumn{4}{r}{Continued on next page} \\
\midrule
\endfoot
\bottomrule
\endlastfoot
Primary & PI (bits/event) & 0.162 [0.086, 0.217] & 9982/10000 \\
Primary & Phase (degrees) & 330.043 [315.889, 343.663] & 9982/10000 \\
Primary & Rate ratio & 4.526 [2.922, 5.958] & 9982/10000 \\
Barker & PI (bits/event) & 0.078 [0.053, 0.106] & 10000/10000 \\
Barker & Phase (degrees) & 340.463 [331.274, 351.668] & 10000/10000 \\
Barker & Rate ratio & 2.657 [2.234, 3.177] & 10000/10000 \\
\end{longtable}
}
\Needspace{8\baselineskip}
\subsection*{(d) Primary effect precision}
{\small
\begin{longtable}{llrrrr}
\toprule
Scenario & Quantity & Center & Lower & Upper & Fits \\
\midrule
\endfirsthead
\toprule
Scenario & Quantity & Center & Lower & Upper & Fits \\
\midrule
\endhead
\midrule
\multicolumn{6}{r}{Continued on next page} \\
\midrule
\endfoot
\bottomrule
\endlastfoot
A: chronology & Rate ratio & 4.53 & 2.92 & 5.96 & 9982 \\
B: sampling & Rate ratio & 4.94 & 1.56 & 15.58 & 5000 \\
C: joint & Rate ratio & 5.09 & 2.10 & 14.19 & 10000 \\
A: chronology & Phase (deg) & 330.04 & 315.89 & 343.66 & 9982 \\
B: sampling & Phase (deg) & 330.00 & 286.53 & 374.08 & 5000 \\
C: joint & Phase (deg) & 329.22 & 295.92 & 365.45 & 10000 \\
\end{longtable}
}
A/C are central 95\% working ranges centered on medians. B projects an approximate joint 95\% coefficient confidence region and is centered on the point estimate. Angles are unwrapped around 330 degrees. C is not a calibrated combined interval; no Barker sampling interval is available (Text S5).

\clearpage
\setcounter{table}{3}\refstepcounter{table}\label{tab:S4}
\section*{Table S4. Design and chronology sensitivity}
% Generated by paper_tables.py; edit the exporter, not this file.
\Needspace{8\baselineskip}
\subsection*{(a) All 30 designs on 173-kyr common support; 55 events}
{\small
\begin{longtable}{rrrrrrrr}
\toprule
History & Bin & Origin & PI & Nominal p & Phase & Rate ratio & dAICc \\
\midrule
\endfirsthead
\toprule
History & Bin & Origin & PI & Nominal p & Phase & Rate ratio & dAICc \\
\midrule
\endhead
\midrule
\multicolumn{8}{r}{Continued on next page} \\
\midrule
\endfoot
\bottomrule
\endlastfoot
1 & 0.1 & 0 & 0.1761 & 0.00121 & 330.7 & 4.74 & -9.40 \\
1 & 0.1 & 0.5 & 0.1799 & 0.00105 & 330.7 & 4.82 & -9.68 \\
1 & 0.2 & 0 & 0.1595 & 0.00228 & 331.1 & 4.37 & -8.10 \\
1 & 0.2 & 0.5 & 0.1620 & 0.00207 & 331.4 & 4.42 & -8.29 \\
1 & 0.5 & 0 & 0.1470 & 0.00368 & 330.9 & 4.08 & -7.05 \\
1 & 0.5 & 0.5 & 0.1666 & 0.00175 & 330.5 & 4.52 & -8.55 \\
1.5 & 0.1 & 0 & 0.1818 & 0.00098 & 329.8 & 5.02 & -9.83 \\
1.5 & 0.1 & 0.5 & 0.1863 & 0.00082 & 329.8 & 5.13 & -10.18 \\
1.5 & 0.2 & 0 & 0.1725 & 0.00139 & 330.0 & 4.76 & -9.09 \\
1.5 & 0.2 & 0.5 & 0.1681 & 0.00165 & 330.6 & 4.66 & -8.76 \\
1.5 & 0.5 & 0 & 0.1596 & 0.00228 & 329.6 & 4.48 & -8.01 \\
1.5 & 0.5 & 0.5 & 0.1725 & 0.00139 & 329.5 & 4.80 & -9.00 \\
2 & 0.1 & 0 & 0.1723 & 0.00141 & 329.3 & 4.92 & -9.10 \\
2 & 0.1 & 0.5 & 0.1835 & 0.00092 & 329.0 & 5.19 & -9.96 \\
2 & 0.2 & 0 & 0.1705 & 0.00150 & 329.1 & 4.87 & -8.94 \\
2 & 0.2 & 0.5 & 0.1729 & 0.00137 & 329.4 & 4.92 & -9.12 \\
2 & 0.5 & 0 & 0.1638 & 0.00194 & 328.5 & 4.66 & -8.33 \\
2 & 0.5 & 0.5 & 0.1499 & 0.00330 & 330.0 & 4.40 & -7.27 \\
3 & 0.1 & 0 & 0.1852 & 0.00086 & 326.6 & 5.41 & -10.09 \\
3 & 0.1 & 0.5 & 0.1891 & 0.00074 & 326.7 & 5.51 & -10.39 \\
3 & 0.2 & 0 & 0.1769 & 0.00118 & 326.8 & 5.19 & -9.43 \\
3 & 0.2 & 0.5 & 0.1786 & 0.00110 & 327.1 & 5.22 & -9.56 \\
3 & 0.5 & 0 & 0.1570 & 0.00251 & 327.1 & 4.65 & -7.82 \\
3 & 0.5 & 0.5 & 0.1677 & 0.00167 & 327.1 & 4.92 & -8.63 \\
5 & 0.1 & 0 & 0.1584 & 0.00239 & 324.5 & 4.86 & -8.04 \\
5 & 0.1 & 0.5 & 0.1614 & 0.00213 & 324.5 & 4.94 & -8.27 \\
5 & 0.2 & 0 & 0.1513 & 0.00312 & 324.9 & 4.67 & -7.48 \\
5 & 0.2 & 0.5 & 0.1480 & 0.00354 & 326.0 & 4.57 & -7.23 \\
5 & 0.5 & 0 & 0.1301 & 0.00702 & 327.7 & 4.06 & -5.76 \\
5 & 0.5 & 0.5 & 0.1602 & 0.00222 & 323.8 & 4.85 & -8.06 \\
\end{longtable}
}
History and bin width are in kyr; origin is a fraction of bin width. Phase is in degrees and PI in bits/event throughout Tables S4--S5.
\Needspace{8\baselineskip}
\subsection*{(b) Background shape and memory alternatives}
{\small
\begin{longtable}{P{3.3cm}rrrrrr}
\toprule
Specification & kyr & PI & Nominal p & Phase & Rate ratio & dAICc \\
\midrule
\endfirsthead
\toprule
Specification & kyr & PI & Nominal p & Phase & Rate ratio & dAICc \\
\midrule
\endhead
\midrule
\multicolumn{7}{r}{Continued on next page} \\
\midrule
\endfoot
\bottomrule
\endlastfoot
Linear climate & 180 & 0.1793 & 0.00107 & 330.0 & 4.94 & -9.61 \\
+ LR04 squared & 180 & 0.1705 & 0.00150 & 322.2 & 4.70 & -8.93 \\
+ CO$_2$ squared & 180 & 0.1711 & 0.00147 & 329.0 & 4.81 & -8.98 \\
+ both squares & 180 & 0.1798 & 0.00105 & 322.6 & 5.00 & -9.63 \\
+ squares + interaction & 180 & 0.1791 & 0.00108 & 322.8 & 5.01 & -9.57 \\
Count in previous 1.5 kyr & 164.8 & 0.1961 & 0.00074 & 333.2 & 5.28 & -10.35 \\
Elapsed time & 164.8 & 0.1883 & 0.00099 & 332.4 & 5.10 & -9.77 \\
Log elapsed time & 164.8 & 0.1920 & 0.00087 & 332.7 & 5.23 & -10.04 \\
\end{longtable}
}
The five climate specifications have 55 events and 180 kyr; the three memory specifications have 53 events and 164.8 kyr. Each dAICc compares full with its own reduced model on identical support.
\Needspace{8\baselineskip}
\subsection*{(c) Response-support comparison at primary settings}
{\small
\begin{longtable}{rlllll}
\toprule
Exposure (kyr) & PI & Nominal p & Phase & Rate ratio & dAICc \\
\midrule
\endfirsthead
\toprule
Exposure (kyr) & PI & Nominal p & Phase & Rate ratio & dAICc \\
\midrule
\endhead
\midrule
\multicolumn{6}{r}{Continued on next page} \\
\midrule
\endfoot
\bottomrule
\endlastfoot
173 & 0.1725 & 0.00139 & 330.0 & 4.76 & -9.09 \\
180 & 0.1793 & 0.00107 & 330.0 & 4.94 & -9.61 \\
\end{longtable}
}
\Needspace{8\baselineskip}
\subsection*{(d) Segment-specific phase coefficients versus a common response}
{\small
\begin{longtable}{rllll}
\toprule
Events & LR & Nominal p & PI & dAICc \\
\midrule
\endfirsthead
\toprule
Events & LR & Nominal p & PI & dAICc \\
\midrule
\endhead
\midrule
\multicolumn{5}{r}{Continued on next page} \\
\midrule
\endfoot
\bottomrule
\endlastfoot
55 & 0.4276 & 0.80752 & 0.00561 & 3.649 \\
\end{longtable}
}
Both phase interactions are added together (2 degrees of freedom); response exposure is 180 kyr. This test does not establish equivalence between segments.
\Needspace{8\baselineskip}
\subsection*{(e) Alternative chronology models}
{\small
\begin{longtable}{lrrr}
\toprule
Case & Valid/total & PI median [95\% range] & p < 0.05 (\%) \\
\midrule
\endfirsthead
\toprule
Case & Valid/total & PI median [95\% range] & p < 0.05 (\%) \\
\midrule
\endhead
\midrule
\multicolumn{4}{r}{Continued on next page} \\
\midrule
\endfoot
\bottomrule
\endlastfoot
NGRIP knots 2.5 kyr & 1993/2000 & 0.161 [0.085, 0.218] & 98.29 \\
NGRIP knots 5 kyr & 1997/2000 & 0.161 [0.085, 0.218] & 98.40 \\
NGRIP knots 10 kyr & 1997/2000 & 0.162 [0.089, 0.213] & 99.00 \\
Sofular So-57 overlap & 9982/10000 & 0.161 [0.086, 0.216] & 98.53 \\
\end{longtable}
}
Percentages are conditional robustness fractions, not bootstrap p values. The So-57 case changes only four supported Sofular projections; the oldest event retains So-4.
\Needspace{8\baselineskip}
\subsection*{(f) Barker event definitions}
{\small
\begin{longtable}{lrrrrrr}
\toprule
Definition & Events & PI & Nominal p & Phase & Rate ratio & dAICc \\
\midrule
\endfirsthead
\toprule
Definition & Events & PI & Nominal p & Phase & Rate ratio & dAICc \\
\midrule
\endhead
\midrule
\multicolumn{7}{r}{Continued on next page} \\
\midrule
\endfoot
\bottomrule
\endlastfoot
variable threshold & 70 & 0.0894 & 0.01309 & 336.9 & 2.88 & -4.65 \\
fixed threshold & 59 & 0.1042 & 0.01409 & 350.3 & 3.10 & -4.50 \\
\end{longtable}
}
\Needspace{8\baselineskip}
\subsection*{(g) Sofular projection-position diagnostic}
{\small
\begin{longtable}{lrrr}
\toprule
Component & Coordinate shift (kyr) & Minimum SD (kyr) & Maximum SD (kyr) \\
\midrule
\endfirsthead
\toprule
Component & Coordinate shift (kyr) & Minimum SD (kyr) & Maximum SD (kyr) \\
\midrule
\endhead
\midrule
\multicolumn{4}{r}{Continued on next page} \\
\midrule
\endfoot
\bottomrule
\endlastfoot
So-4 & -1 & 0.259 & 0.429 \\
So-4 & -0.5 & 0.289 & 0.422 \\
So-4 & 0 & 0.254 & 0.403 \\
So-4 & 0.5 & 0.254 & 0.385 \\
So-4 & 1 & 0.246 & 0.368 \\
So-57 & -1 & 0.228 & 0.364 \\
So-57 & -0.5 & 0.237 & 0.402 \\
So-57 & 0 & 0.270 & 0.444 \\
So-57 & 0.5 & 0.248 & 0.489 \\
So-57 & 1 & 0.232 & 0.536 \\
\end{longtable}
}
Ranges span supported events and all nine detector settings before order conditioning. Coordinate shifts move error-field evaluation positions without changing nominal event ages; they do not represent climate lags. The complete 405-row diagnostic accompanies this table.

\clearpage
\setcounter{table}{4}\refstepcounter{table}\label{tab:S5}
\section*{Table S5. Orbital and climate--phase comparisons}
% Generated by paper_tables.py; edit the exporter, not this file.
\Needspace{8\baselineskip}
\subsection*{(a) Primary: all point-age comparisons}
{\small
\begin{longtable}{P{5cm}rrrrrr}
\toprule
Comparison & df & PI & Nominal p & Holm p & dAIC & dAICc \\
\midrule
\endfirsthead
\toprule
Comparison & df & PI & Nominal p & Holm p & dAIC & dAICc \\
\midrule
\endhead
\midrule
\multicolumn{7}{r}{Continued on next page} \\
\midrule
\endfoot
\bottomrule
\endlastfoot
BG + Pre vs BG & 2 & 0.1793 & 0.00107 & -- & -9.67 & -9.61 \\
BG + Ecc vs BG & 1 & 0.0389 & 0.08496 & 0.42479 & -0.97 & -0.94 \\
BG + Pre + Ecc vs BG + Pre & 1 & 0.0125 & 0.32914 & 0.95927 & 1.05 & 1.08 \\
BG + Pre + Ecc vs BG + Ecc & 2 & 0.1529 & 0.00294 & 0.02059 & -7.66 & -7.59 \\
BG + Obl vs BG & 1 & 0.0001 & 0.92315 & 0.95927 & 1.99 & 2.02 \\
BG + Pre + Obl vs BG + Pre & 1 & 0.0130 & 0.31976 & 0.95927 & 1.01 & 1.05 \\
BG + Pre + Obl vs BG + Obl & 2 & 0.1922 & 0.00066 & 0.00592 & -10.65 & -10.59 \\
BG + Insol vs BG & 1 & 0.1225 & 0.00224 & 0.01790 & -7.34 & -7.32 \\
BG + Pre + Insol vs BG + Pre & 1 & 0.0219 & 0.19617 & 0.78467 & 0.33 & 0.37 \\
BG + Pre + Insol vs BG + Insol & 2 & 0.0787 & 0.04977 & 0.29865 & -2.00 & -1.93 \\
\end{longtable}
}
\Needspace{8\baselineskip}
\subsection*{(b) Primary: chronology sensitivity}
{\small
\begin{longtable}{P{6cm}rrr}
\toprule
Comparison & PI median [2.5\%, 97.5\%] & Valid & Total \\
\midrule
\endfirsthead
\toprule
Comparison & PI median [2.5\%, 97.5\%] & Valid & Total \\
\midrule
\endhead
\midrule
\multicolumn{4}{r}{Continued on next page} \\
\midrule
\endfoot
\bottomrule
\endlastfoot
BG + Pre vs BG & 0.164 [0.086, 0.217] & 499 & 500 \\
BG + Ecc vs BG & 0.033 [0.008, 0.056] & 499 & 500 \\
BG + Pre + Ecc vs BG + Pre & 0.011 [0.000, 0.027] & 499 & 500 \\
BG + Pre + Ecc vs BG + Ecc & 0.140 [0.077, 0.189] & 499 & 500 \\
BG + Obl vs BG & 0.000 [0.000, 0.005] & 499 & 500 \\
BG + Pre + Obl vs BG + Pre & 0.013 [0.002, 0.034] & 499 & 500 \\
BG + Pre + Obl vs BG + Obl & 0.179 [0.091, 0.236] & 499 & 500 \\
BG + Insol vs BG & 0.108 [0.042, 0.146] & 499 & 500 \\
BG + Pre + Insol vs BG + Pre & 0.019 [0.001, 0.042] & 499 & 500 \\
BG + Pre + Insol vs BG + Insol & 0.072 [0.040, 0.127] & 499 & 500 \\
\end{longtable}
}
\Needspace{8\baselineskip}
\subsection*{(c) Barker: all point-age comparisons}
{\small
\begin{longtable}{P{5cm}rrrrrr}
\toprule
Comparison & df & PI & Nominal p & Holm p & dAIC & dAICc \\
\midrule
\endfirsthead
\toprule
Comparison & df & PI & Nominal p & Holm p & dAIC & dAICc \\
\midrule
\endhead
\midrule
\multicolumn{7}{r}{Continued on next page} \\
\midrule
\endfoot
\bottomrule
\endlastfoot
BG + Pre vs BG & 2 & 0.0894 & 0.01309 & -- & -4.67 & -4.65 \\
BG + Ecc vs BG & 1 & 0.0110 & 0.30211 & 1.00000 & 0.94 & 0.95 \\
BG + Pre + Ecc vs BG + Pre & 1 & 0.0050 & 0.48443 & 1.00000 & 1.51 & 1.53 \\
BG + Pre + Ecc vs BG + Ecc & 2 & 0.0834 & 0.01746 & 0.13968 & -4.10 & -4.07 \\
BG + Obl vs BG & 1 & 0.0007 & 0.79845 & 1.00000 & 1.93 & 1.94 \\
BG + Pre + Obl vs BG + Pre & 1 & 0.0024 & 0.62724 & 1.00000 & 1.76 & 1.78 \\
BG + Pre + Obl vs BG + Obl & 2 & 0.0911 & 0.01202 & 0.10819 & -4.84 & -4.82 \\
BG + Insol vs BG & 1 & 0.0521 & 0.02450 & 0.17148 & -3.06 & -3.05 \\
BG + Pre + Insol vs BG + Pre & 1 & 0.0011 & 0.74646 & 1.00000 & 1.90 & 1.91 \\
BG + Pre + Insol vs BG + Insol & 2 & 0.0383 & 0.15590 & 0.93543 & 0.28 & 0.31 \\
\end{longtable}
}
\Needspace{8\baselineskip}
\subsection*{(d) Barker: chronology sensitivity}
{\small
\begin{longtable}{P{6cm}rrr}
\toprule
Comparison & PI median [2.5\%, 97.5\%] & Valid & Total \\
\midrule
\endfirsthead
\toprule
Comparison & PI median [2.5\%, 97.5\%] & Valid & Total \\
\midrule
\endhead
\midrule
\multicolumn{4}{r}{Continued on next page} \\
\midrule
\endfoot
\bottomrule
\endlastfoot
BG + Pre vs BG & 0.078 [0.054, 0.106] & 500 & 500 \\
BG + Ecc vs BG & 0.010 [0.008, 0.013] & 500 & 500 \\
BG + Pre + Ecc vs BG + Pre & 0.005 [0.003, 0.009] & 500 & 500 \\
BG + Pre + Ecc vs BG + Ecc & 0.073 [0.050, 0.100] & 500 & 500 \\
BG + Obl vs BG & 0.000 [0.000, 0.002] & 500 & 500 \\
BG + Pre + Obl vs BG + Pre & 0.002 [0.000, 0.005] & 500 & 500 \\
BG + Pre + Obl vs BG + Obl & 0.079 [0.055, 0.108] & 500 & 500 \\
BG + Insol vs BG & 0.048 [0.030, 0.069] & 500 & 500 \\
BG + Pre + Insol vs BG + Pre & 0.001 [0.000, 0.004] & 500 & 500 \\
BG + Pre + Insol vs BG + Insol & 0.031 [0.018, 0.049] & 500 & 500 \\
\end{longtable}
}
\Needspace{8\baselineskip}
\subsection*{(e) Additional LR04-by-phase interactions (2 degrees of freedom)}
{\small
\begin{longtable}{llllllr}
\toprule
Catalogue & PI & LR & Nominal p & dAIC & dAICc & Valid MC \\
\midrule
\endfirsthead
\toprule
Catalogue & PI & LR & Nominal p & dAIC & dAICc & Valid MC \\
\midrule
\endhead
\midrule
\multicolumn{7}{r}{Continued on next page} \\
\midrule
\endfoot
\bottomrule
\endlastfoot
Primary & 0.0258 & 1.9709 & 0.37327 & 2.03 & 2.11 & 499 \\
Barker & 0.0561 & 5.4464 & 0.06566 & -1.45 & -1.42 & 500 \\
\end{longtable}
}
BG denotes intercept, prior-event count, LR04 and CO2, plus a segment term in the primary catalogue; Pre is the sine/cosine phase block; Ecc is eccentricity; Obl is obliquity; Insol is 65-degree-N summer-solstice daily mean insolation. Holm adjustment includes nine additional orbital comparisons per catalogue; the phase reference is outside that family. LR04 interactions are separate exploratory tests. All p values in this table are nominal; the point-age main null bootstrap is not reused for different hypotheses. Negative dAIC/dAICc favors the augmented model.

\clearpage
\begin{figure}[p]
\centering\paperfigure{FigS01.pdf}
\caption{Sofular and MF proxy correspondence across their overlapping interval. (a) Sofular $\delta^{13}$C, with labels from \citet{held2024dansgaard}. (b) MF $\delta^{18}$O, with labels from \citet{fohlmeister2023role}. Small points and fine lines show observations; thicker curves show 50-yr Gaussian smoothing after interpolation to a 1-yr grid, independently within segments separated by gaps exceeding 0.5 kyr. Gray connectors identify six visually inferred feature correspondences; they guide comparison and do not demonstrate exact synchronization. MF 6.16 corresponds to an unnumbered Sofular feature near 175.2 kyr BP, whereas Sofular 6.9, near 176.8 kyr BP, is the next older distinct published event retained in the composite. Fine dashed leaders terminate at approximate label anchors rather than precise transition ages. Gray shading denotes the 7.08-kyr MF proxy gap. The Sofular proxy axis is reversed. Ages are kyr BP relative to 1950 CE.}
\label{fig:S01}
\end{figure}
\clearpage
\begin{figure}[p]
\centering\paperfigure{FigS02.pdf}
\caption{The 21-event MF--Sofular sequence and numerical timing assignments. (a) Ordered composite with source records; the dashed line at 175 kyr BP indicates the displayed MF--Sofular division. (b) MF $\delta^{18}$O and published events 6.1--6.16 \citep{fohlmeister2023role}. (c) Sofular $\delta^{13}$C; source labels 6.9--6.13 are shown in brackets and renumbered 6.17--6.21 \citep{held2024dansgaard}. Gray points and colored curves show observations and Gaussian-smoothed records. Filled circles mark the largest expected-direction gradient within $\pm400$ yr of each literature anchor after gap-bounded 1-yr interpolation and 50-yr Gaussian smoothing; leaders connect labels to assigned times. MF 6.16 and Sofular 6.9 are retained as distinct events following Figure~\ref{fig:S01}. Open gray squares and horizontal bars near the bottom of (c) show So-4/So-57 U--Th control ages and reported $2\sigma$ errors, not event-age intervals. Controls outside the plotting window remain in the processed source table. The Sofular axis is reversed; ages are kyr BP relative to 1950 CE. ``MIS6'' labels denote the project sequence, whose support extends beyond strict MIS~6.}
\label{fig:S02}
\end{figure}
\clearpage
\begin{figure}[p]
\centering\paperfigure{FigS03.pdf}
\caption{NGRIP working event-age uncertainty. (a) Input error scales for 34 warming (GI; orange circles) and 35 cooling (GS; blue squares) boundaries from \citet{rasmussen2014stratigraphic}. Symbols show adopted definition-error standard deviations; displaced event labels have thin connectors to their values. The black line connects event-level chronology envelopes: published maximum counting errors in the counted section and the study's nominal $\pm4.5\%$-of-b2k-age envelope in the model extension. Chronology envelopes are interpreted approximately as $2\sigma$; definition errors use working $1\sigma$ values. The vertical scale is logarithmic. (b) Combined chronology-plus-definition offsets from point ages. Gray lines show 30 evenly indexed realizations; blue line and shading show the median and pointwise 2.5th--97.5th percentiles of all 10,000 jointly ordered realizations. The dashed boundary is 60.202 ka b2k, or 60.152 kyr BP. Positive offsets indicate older ages. The model-extension envelope is a study scenario informed by \citet{moseley2020nalps19}, not a published pointwise uncertainty curve. Gaussian envelopes are not hard bounds; MC quantiles describe the adopted model rather than calibrated confidence intervals.}
\label{fig:S03}
\end{figure}
\clearpage
\begin{figure}[p]
\centering\paperfigure{FigS04.pdf}
\caption{Chronology and event-timing sensitivity of the speleothem composite. (a) Sampled age offsets from nominal maximum-gradient picks, with nominal ages beside event labels. Points mark medians; thick and thin intervals span the 16th--84th and 2.5th--97.5th percentiles. The dashed zero line denotes the nominal sequence, and colors distinguish MF and Sofular. (b) Pearson correlations among sampled age anomalies; the white boundary separates records. The 10,000 realizations choose one of nine smoothing/window settings per record and preserve event identity and order. MF shares a Gaussian chronology factor, scaled by its repaired approximate 95\% age envelope. Sofular uses independent Gaussian U--Th errors at ten So-4 dated depths, with standard deviations equal to reported $2\sigma$ errors divided by two; a monotone warp preserves the published reference curve between dates. Nonmonotone knots and crossed complete event sequences are rejected. No interpolation crosses a growth hiatus or extrapolates outside support. Projecting stack ages onto an individual curve locates an assumed error-field coordinate, not a verified event layer. The ensemble does not reproduce StalAge/iscam posterior ages or include proxy noise, unknown growth history, shared systematic dating errors, or stack-alignment lags. Ages are kyr BP relative to 1950 CE.}
\label{fig:S04}
\end{figure}
\clearpage
\begin{figure}[p]
\centering\paperfigure{FigS05.pdf}
\caption{Barker SpeleoAge control errors and chronology realizations \citep{barker2011800}. (a) Printed combined matching and absolute-dating uncertainties for 60 Supplementary Table~S1 controls (black circles), used as uniform proposal half-widths. The line joins these half-widths and represents an interpolated offset bound, not a standard deviation. The blue diamond is an auxiliary 400-kyr control with a 1.30-kyr half-width, chosen as the larger of the last-four-control linear extrapolation and their maximum error. (b) Twenty-five accepted age-offset maps (gray), median (blue), and pointwise 2.5th--97.5th percentiles (shading) of 10,000 accepted maps. Positive offsets mean older ages. Entire crossed control proposals are rejected, so accepted medians need not equal zero. Red ticks locate the 70 nominal events. Ochre shading marks the approximately 265--315-kyr alignment gap; offsets interpolate between its bracketing controls at 264.24 and 317.70 kyr. The blue dotted line identifies the final published control at 378.15 kyr and pale blue shading its extension to the auxiliary endpoint. All weights use SpeleoAge. BP1950 remains the working epoch assumption. Ranges do not provide calibrated confidence intervals and include no additional event-picking or membership uncertainty.}
\label{fig:S05}
\end{figure}
\clearpage
\begin{figure}[p]
\centering\paperfigure{FigS06.pdf}
\caption{Primary conditional results under chronology and transition-timing perturbations. Histograms summarize 9,982 supported fits from 10,000 independently paired source realizations. Black solid lines show point-age estimates, red dashed lines MC medians, and gray dotted lines decision thresholds. (a) PI has median 0.16215 bits/event and central 95\% range 0.08616--0.21746, versus 0.17933 at point ages. (b) Median nominal likelihood-ratio p is 0.002067 (0.000251--0.03780); 9,851/9,982 fits have nominal $p<0.05$. (c) Preferred phase is $330.04^\circ$ ($315.89$--$343.66^\circ$). (d) Maximum/minimum phase rate ratio is 4.53 (2.92--5.96). (e) Full-minus-reduced AICc is $-8.30$ ($-12.52$ to $-2.49$), negative in 9,971 fits. Every realization updates counts and 1.5-kyr history, retaining the original 0.2-kyr bins, forcing grids, scaling, and 180-kyr response exposure. Eighteen draws leave observation support and remain recorded without PI estimates; 66 valid draws have 54 response events and 9,916 have 55. Ages are neither clipped nor redrawn to preserve response membership. Quantiles and proportions are conditional on supported fits and describe the working age model, not complete sampling uncertainty. The 98.69\% threshold fraction is not an empirical p value; the known fraction across all 10,000 draws is at least 98.51\%.}
\label{fig:S06}
\end{figure}
\clearpage
\begin{figure}[p]
\centering\paperfigure{FigS07.pdf}
\caption{Barker precession association under chronology perturbations. All 10,000 ordered realizations retain the same 70 variable-threshold events. Histograms show (a) conditional PI, (b) nominal likelihood-ratio p, (c) preferred phase, (d) maximum/minimum phase rate ratio, and (e) full-minus-reduced AICc. Black solid and red dashed lines mark point-age estimates and MC medians; gray dotted lines denote $p=0.05$ and zero AICc difference. Preferred phases are unwrapped around the point estimate, with axis labels modulo $360^\circ$. (f) Fitted phase multipliers: point ages (black), pointwise MC median (red dashed), and 2.5th--97.5th percentiles (blue shading); the horizontal reference is unity. Phase 0 denotes a precession-index minimum. Every realization uses 0.2-kyr nominal bins and recomputes the preceding 1.5-kyr event count. Reduced/full models retain fixed LR04, \ce{CO2}, phase grids, and scaling; response support is 0--398.5 kyr, with 398.5--400 kyr used for history only. PI median is 0.07763 bits/event (0.05350--0.10610); 89.66\% of nominal p values are below 0.05. These are age-sensitivity ranges and a robustness fraction, not event-process confidence intervals or an empirical p value.}
\label{fig:S07}
\end{figure}
\clearpage
\begin{figure}[p]
\centering\paperfigure{FigS08.pdf}
\caption{Primary reduced-model bootstrap calibration. Gray bars show the likelihood-ratio distribution from 9,999 catalogues generated under the fitted reduced Poisson model; the blue dashed curve is the nominal $\chi^2_2$ reference, and the magenta line marks observed LR 13.6732. Simulation proceeds from older to younger bins, dynamically updating the preceding 1.5-kyr event count, with separate zero-history initialization at each observation boundary. The history buffers are simulated, and history resets across the unobserved gap. Both reduced and phase-augmented models are refitted. Twenty-one null replicates equal or exceed the observed LR, giving $(21+1)/(9999+1)=0.0022$. The exact binomial 95\% interval for the underlying unadjusted exceedance probability is 0.001301--0.003209. All replicates pass convergence, nesting, and clipping checks. Event totals vary; fixed forcing grids and point-age fits are used. Chronology and transition-timing MC are separate from this null calibration.}
\label{fig:S08}
\end{figure}
\clearpage
\begin{figure}[p]
\centering\paperfigure{FigS09.pdf}
\caption{Barker reduced-model bootstrap calibration for the primary variable-threshold catalogue. Gray bars show the density of 9,999 simulated LR statistics. The reduced model includes an intercept, preceding 1.5-kyr event count, LR04, and \ce{CO2}; full adds phase sine and cosine. Histories update from older to younger bins, and both models are refitted for each simulated catalogue. The colored line marks observed LR 8.6716; the dashed curve is the nominal $\chi^2_2$ reference. There are 137 exceedances, giving plus-one bootstrap $p=0.0138$. The displayed exact binomial interval, 0.01152--0.01618, describes finite-simulation precision of the unadjusted null exceedance probability, not an age interval or effect confidence interval. Simulation retains 0--400-kyr observation and 0--398.5-kyr response support and allows event totals to vary. All 9,999 replicates are valid without retries. Fixed-threshold events and chronology MC are not included.}
\label{fig:S09}
\end{figure}
\clearpage
\begin{figure}[p]
\centering\paperfigure{FigS10.pdf}
\caption{Primary sensitivity to analysis design and response support. Thirty designs cross five history windows (1, 1.5, 2, 3, and 5 kyr), three bin widths (0.1, 0.2, and 0.5 kyr; colors), and two origins (unshifted: filled circles/solid lines; half-bin shifted: open squares/dashed lines). All use the 173-kyr common response core, NGRIP 12--118 and speleothems 132.5--199.5 kyr BP. The star marks the primary 1.5-kyr history, 0.2-kyr unshifted design; this history window was motivated by D--O recurrence, not chosen by the p values shown. (a) PI spans 0.13009--0.18908 bits/event. (b) Nominal LR p spans 0.000740--0.007016; the dotted line is 0.05. (c) Preferred phase spans $323.82$--$331.39^\circ$. Full-minus-reduced AICc is negative in all designs. (d) At primary settings, the common core gives PI 0.17249 and nominal p 0.001394, whereas 180-kyr main support gives 0.17933 and 0.001074. All designs retain 55 events and exclude the inter-segment gap from both exposure and history. No alternative design inherits the main null-bootstrap calibration.}
\label{fig:S10}
\end{figure}
\clearpage
\begin{figure}[p]
\centering\paperfigure{FigS11.pdf}
\caption{Primary sensitivity to background-climate shape and event memory. All fits use point ages, 0.2-kyr bins, and primary climate scaling; the target is warming-onset rate per unit of whole observation time. (a--c) PI, preferred phase, and maximum/minimum phase rate ratio for five background specifications on identical 180-kyr support and 55 events: linear LR04/\ce{CO2}, plus LR04 squared, \ce{CO2} squared, both squares, or both squares and their product. Each specification enters both reduced and full models, with only phase sine/cosine added in full. (d--f) Corresponding quantities for preceding 1.5-kyr count, elapsed time since the last event, and $\log(1+\mathrm{elapsed\ time}/1\ {\rm kyr})$. These models share 164.8-kyr support and 53 events. Each segment's oldest occupied bin initializes history but contributes no response; its and older exposure are excluded. Elapsed time uses strictly older occupied bin centers and never crosses the gap. Younger event-free tails remain exposed. Points are fitted estimates without confidence intervals, and dashed vertical lines identify reference models on each row's support. All eight variants have positive phase likelihood gain and negative within-variant full-minus-reduced AICc. The two rows have different support and cannot be compared as if only memory changed.}
\label{fig:S11}
\end{figure}
\clearpage
\begin{figure}[p]
\centering\paperfigure{FigS12.pdf}
\caption{Barker event-definition sensitivity on SpeleoAge \citep{barker2011800}. Rose denotes 70 variable-threshold events (primary); blue denotes 59 fixed-threshold events (sensitivity). (a) Events on the La2004 precession index, using filled circles and open squares; shared picks display both symbols. The gray 398.5--400-kyr interval supplies history only, and all point-age events lie in 0--398.5-kyr response support. BP1950 is assumed for SpeleoAge; La2004 receives its documented J2000 correction. (b) Event counts in twelve $30^\circ$ phase sectors, with filled rose bars and dashed blue outlines. Phase increases toward older ages, from minima ($0^\circ$) through maxima ($180^\circ$). Arrow directions are circular means, and lengths equal the mean resultant length times the largest sector count across catalogues; radial ticks are event counts. Text shows mean resultant length and nominal Rayleigh p. (c) Fitted phase multipliers (rose solid; blue dashed); dotted vertical lines mark preferred phases and the horizontal reference is unity. Models share support and settings but recompute histories for each definition. Preferred phases are $336.9^\circ$ and $350.3^\circ$, with rate ratios 2.88 and 3.10. PI is fitted likelihood gain in bits/event; panel p values are nominal $\chi^2_2$ tests. These are alternative definitions of one record, not independent samples. Bootstrap and age-MC results apply only to the variable-threshold catalogue.}
\label{fig:S12}
\end{figure}
\clearpage
\begin{figure}[p]
\centering\paperfigure{FigS13.pdf}
\caption{Exploratory LR04 modulation of the primary phase response. (a) Phase multipliers from a continuous model with LR04-by-sine and LR04-by-cosine terms, evaluated at exposure-weighted LR04 quartiles (legend values in per mil). Phase is zero at precession minima and $180^\circ$ at maxima, increasing toward older ages. Lines use point ages; shading shows pointwise 2.5th--97.5th percentiles from chronology MC. Multipliers are relative to the same background/history with the phase contribution set to zero; they are not total warming rates. The dashed horizontal line denotes unity. (b) Additional PI of the two interaction terms relative to the full model, which includes history, LR04, \ce{CO2}, a segment indicator, and phase sine/cosine. Solid and dashed vertical lines show point-age PI and MC median; the pale band spans the central 95\% of valid chronology results. Within each draw, compared models share events, bins, and exposure. Of 500 selected chronologies, 499 are supported. Point LR is 1.9709 (2 df), with nominal $p=0.3733$. The interaction lacks a separate null bootstrap; the chronology ranges are not complete sampling intervals, and separated curves do not establish modulation.}
\label{fig:S13}
\end{figure}
\clearpage
\begin{figure}[p]
\centering\paperfigure{FigS14.pdf}
\caption{Exploratory LR04 modulation of the Barker variable-threshold phase response. (a) Phase multipliers from the model containing LR04-by-sine and LR04-by-cosine terms, evaluated at exposure-weighted LR04 quartiles (per mil). Point-age curves and pointwise 2.5th--97.5th-percentile age bands show the phase contribution at a fixed background/history, relative to a phase contribution of zero. Phase convention is identical to Figure~\ref{fig:S13}; the horizontal dashed reference is unity. (b) Additional PI of the two interaction terms relative to full, which retains intercept, history, LR04, \ce{CO2}, and both phase terms. Solid and dashed vertical lines denote the point result and MC median; the pale band spans the central 95\% age-sensitivity range. All 500 selected chronologies yield valid comparisons on identical within-draw events, bins, and exposure. Point LR is 5.4464 (2 df), nominal $p=0.06566$, without a separate null bootstrap. Positive nested-model PI and separated fitted curves do not establish modulation; age sensitivity does not include complete event-process sampling uncertainty. The fixed-threshold definition is not included.}
\label{fig:S14}
\end{figure}
\clearpage
\bibliographystyle{agu}
\bibliography{reference}
\end{document}
